% ============================================================
%  CHAPITRE 2 — Analyse et spécification des besoins
% ============================================================

\chapter{Analyse et spécification des besoins}

% ------------------------------------------------------------
\section*{Introduction}
% ------------------------------------------------------------

Avant d'écrire la première ligne de code, toute démarche sérieuse
de développement passe par une phase d'analyse. Ce chapitre en
constitue le compte-rendu : nous y formalisons ce que la plateforme
doit faire, pour qui elle doit le faire, et dans quelles conditions
elle doit fonctionner.

Nous débutons par les besoins d'affaires, qui donnent les raisons
profondes du projet et servent de boussole pour tous les arbitrages
qui suivront. Nous en déduisons ensuite les exigences fonctionnelles
et non fonctionnelles du système, en identifiant les parties prenantes
et en cartographiant les interactions entre acteurs et modules.

La planification du développement, organisée en cinq sprints Agile,
fait l'objet d'une section dédiée. Nous présentons ensuite
l'architecture générale retenue, puis l'environnement technique
mobilisé. Enfin, une inspection approfondie des données sources
débouche sur la conception du modèle physique de l'entrepôt,
sur lequel reposent les tableaux de bord et les modules d'IA.

% ============================================================
\section{Définition des besoins d'affaires}
% ============================================================

Le diagnostic posé au chapitre précédent a mis en évidence les
limites d'une organisation fondée sur des fichiers Excel dispersés.
Pour y répondre, Wevioo a exprimé sept besoins d'affaires
structurants, dont l'objectif commun est de doter l'organisation
d'une plateforme unifiée couvrant l'intégralité du cycle de vie
des projets environnementaux tout en générant une valeur analytique
et décisionnelle réelle.

\begin{itemize}
    \item \textbf{Centralisation des données :} remplacer les fichiers
    Excel dispersés par un référentiel unique consolidant toutes
    les informations projet.

    \item \textbf{Qualité et fiabilité des données :} mettre en place
    des mécanismes de validation à la saisie et un pipeline de
    nettoyage automatisé garantissant la cohérence des données.

    \item \textbf{Suivi opérationnel en temps réel :} offrir aux
    responsables une vision instantanée et actualisée de l'état
    de leur portefeuille.

    \item \textbf{Aide à la décision :} mettre à disposition des
    décideurs des indicateurs agrégés et des tableaux de bord
    interactifs pour l'analyse transversale des projets.

    \item \textbf{Anticipation des risques :} exploiter des modèles
    d'apprentissage automatique pour identifier précocement les
    projets en danger et orienter les arbitrages correctifs.

    \item \textbf{Assistance utilisateur :} intégrer un chatbot
    intelligent permettant d'accéder rapidement à l'information
    et d'être guidé dans les interactions avec la plateforme.

    \item \textbf{Automatisation des traitements :} réduire la charge
    manuelle en automatisant les flux de données et les transformations
    ETL, limitant ainsi les erreurs humaines et les délais de
    traitement.
\end{itemize}

Ces sept axes constituent le socle fonctionnel sur lequel repose
l'ensemble des choix de conception présentés dans ce rapport.

% ============================================================
\section{Spécification des besoins}
% ============================================================

\subsection{Identification des parties prenantes}

Deux profils d'utilisateurs ont été identifiés comme parties
prenantes directes du système.

\begin{itemize}
    \item \textbf{Les administrateurs :} ils assurent la configuration
    globale de la plateforme, la gestion des comptes, le paramétrage
    des formulaires et la supervision des données projet. Leurs
    droits d'accès sont complets.

    \item \textbf{Les agents métiers :} utilisateurs opérationnels
    au quotidien, responsables de l'enregistrement et du suivi des
    projets. Leur périmètre d'accès est limité aux fonctionnalités
    relevant de leur rôle.
\end{itemize}

% ------------------------------------------------------------
\subsection{Besoins fonctionnels}
% ------------------------------------------------------------

Les besoins fonctionnels décrivent ce que le système doit être
capable d'accomplir pour servir les objectifs métiers définis
précédemment. Le tableau~\ref{tab:besoins_fonctionnels} en présente
une synthèse par domaine.

\begin{longtable}{|m{3.5cm}|m{9.5cm}|}
\hline
\textbf{Domaine} & \textbf{Fonctionnalité attendue} \\
\hline
\endhead

Gestion des projets
& Enregistrer, consulter, modifier et supprimer les projets
environnementaux (bases PGD et PPP) via une interface centralisée. \\
\hline

Paramétrage
& Configurer dynamiquement les champs et les formulaires de saisie
selon les besoins spécifiques de chaque type de projet (domaines,
thèmes, catégories, acteurs, financements, données géographiques). \\
\hline

Gestion des accès
& Administrer les comptes utilisateurs, les rôles et les droits
d'accès à la plateforme. \\
\hline

Pipeline ETL
& Collecter, transformer et nettoyer les données issues de la
base Joget (\texttt{jwdb}) afin de garantir leur qualité avant
chargement dans l'entrepôt \texttt{cleaned\_dw}. \\
\hline

Visualisation BI
& Générer des tableaux de bord interactifs avec Apache Superset
et Chart.js permettant le suivi des indicateurs clés et l'analyse
de l'avancement des projets. \\
\hline

Prédiction IA
& Évaluer automatiquement le niveau de risque de chaque projet
à l'aide d'un classificateur XGBClassifier et fournir une
explication contextuelle générée par le modèle de langage
\texttt{llama-3.3-70b-versatile}. \\
\hline

Assistance utilisateur
& Fournir un chatbot RAG capable de répondre aux questions des
utilisateurs sur les projets environnementaux à partir de la base
documentaire interne. \\
\hline

Bibliothèque documentaire
& Permettre l'ajout, la modification, la suppression et le
téléchargement de documents liés aux projets. \\
\hline

\captionsetup{justification=centering}
\caption{Synthèse des besoins fonctionnels}
\label{tab:besoins_fonctionnels}
\end{longtable}

% ------------------------------------------------------------
\subsubsection{Diagramme de cas d'utilisation global}
% ------------------------------------------------------------

La figure~\ref{fig:cas-utilisation-global} offre une vue d'ensemble
des fonctionnalités de la plateforme et illustre la répartition
des responsabilités entre les deux acteurs du système.

\begin{figure}[H]
    \centering
    \includegraphics[width=0.80\textwidth]{img/use_case_general_otpp.png}
    \caption{Diagramme de cas d'utilisation global de la plateforme}
    \label{fig:cas-utilisation-global}
\end{figure}

Ce diagramme met en scène deux acteurs et huit modules fonctionnels.

\paragraph{Acteurs.}
\begin{itemize}
    \item \textbf{L'Administrateur} dispose du périmètre d'action le
    plus étendu : configuration du système, gestion des utilisateurs,
    administration des contenus et accès à l'ensemble des
    fonctionnalités.

    \item \textbf{L'Agent métier} opère dans les limites des
    permissions qui lui sont accordées : consultation, saisie et
    visualisation des tableaux de bord.
\end{itemize}

\paragraph{Modules fonctionnels.}

\begin{enumerate}
    \item \textbf{Tableaux de bord :} les deux acteurs accèdent
    à des dashboards développés avec Apache Superset et Chart.js,
    offrant une visualisation actualisée des indicateurs de
    performance des projets.

    \item \textbf{Bibliothèque des documents :} l'administrateur
    gère le cycle de vie des documents ; l'agent métier peut les
    consulter et les télécharger.

    \item \textbf{Gestion des utilisateurs :} fonctionnalité réservée
    à l'administrateur, couvrant la création, la modification, la
    suppression des comptes et l'assignation des rôles.

    \item \textbf{Base de Gestion des Déchets (PGD) :} gestion
    complète des projets de déchets avec opérations CRUD et
    fonctionnalités de recherche.

    \item \textbf{Base des Projets Publics (PPP) :} consultation et
    gestion des projets publics, avec droits de modification réservés
    à l'administrateur.

    \item \textbf{Diagnostic de risque des projets :} module IA
    central. L'administrateur saisit les indicateurs d'un projet
    et obtient un score de risque accompagné d'une explication
    exportable.

    \item \textbf{Chatbot intelligent :} assistant conversationnel
    accessible aux deux acteurs, fondé sur une architecture RAG
    pour répondre aux questions portant sur les projets
    environnementaux.

    \item \textbf{Modules de paramétrage :} trois sous-modules
    réservés à l'administrateur pour configurer les référentiels :
    \begin{itemize}
        \item Paramétrage PGD (domaines, thèmes, catégories, acteurs,
              financements) ;
        \item Paramétrage PPP (partenariats, types de projets) ;
        \item Paramétrage géographique (gouvernorats, délégations,
              secteurs).
    \end{itemize}
\end{enumerate}

% ------------------------------------------------------------
\subsection{Besoins non fonctionnels}
% ------------------------------------------------------------

Au-delà des fonctionnalités, la plateforme doit satisfaire
six exigences de qualité transversales.

\begin{itemize}
    \item \textbf{Fiabilité :} des résultats cohérents et des données
    exactes sont indispensables pour que les analyses produites et
    les décisions qui en découlent restent crédibles.

    \item \textbf{Sécurité :} la confidentialité des données sensibles
    doit être assurée, et les accès contrôlés via un système de rôles
    et permissions.

    \item \textbf{Disponibilité :} les projets environnementaux
    peuvent nécessiter un suivi à tout moment ; la plateforme doit
    rester accessible en permanence.

    \item \textbf{Performance :} traitements ETL, requêtes analytiques
    et rendu des tableaux de bord doivent s'exécuter dans des délais
    acceptables même sur des volumes de données croissants.

    \item \textbf{Ergonomie :} l'interface doit rester utilisable par
    des profils non techniques, pour favoriser l'adoption par
    l'ensemble des agents métiers.

    \item \textbf{Scalabilité :} l'architecture doit absorber l'ajout
    de nouvelles fonctionnalités et la montée en charge des données
    sans refonte majeure.
\end{itemize}

% ============================================================
\section{Planification du projet}
% ============================================================

Le développement est découpé en cinq sprints Agile. Cette approche
incrémentale livre des fonctionnalités testables à l'issue de
chaque itération, autorise des ajustements de priorité en cours de
route et maintient une communication continue avec les parties
prenantes. Les trois tableaux ci-dessous présentent respectivement
le product backlog, la répartition par release et le calendrier
détaillé.

\vspace{0.4cm}

\begin{longtable}{|m{1.2cm}|m{3.8cm}|m{6.2cm}|m{2cm}|m{1.5cm}|}
\hline
\textbf{ID} & \textbf{Backlog Item} & \textbf{Description} &
\textbf{Priorité} & \textbf{Sprint} \\
\hline
\endhead

B1 & Gestion des projets environnementaux &
Développement des interfaces d'enregistrement et de suivi
des projets PGD et PPP. & Haute & 1 \\
\hline

B2 & Module de paramétrage &
Mise en place du système de paramétrage dynamique des formulaires,
champs et données référentielles. & Haute & 1 \\
\hline

B3 & Gestion des utilisateurs &
Développement du module de gestion des comptes, rôles et
permissions. & Haute & 2 \\
\hline

B4 & Pipeline ETL &
Conception et développement du pipeline Python/Pandas pour le
nettoyage, la transformation et le chargement dans
\texttt{cleaned\_dw}. & Haute & 2 \\
\hline

B5 & Dashboards analytiques &
Développement des tableaux de bord Superset et intégration des
visualisations Chart.js dans Joget. & Moyenne & 3 \\
\hline

B6 & Prédiction des risques &
Conception, entraînement et déploiement du modèle XGBClassifier
avec génération d'explication via LLM. & Haute & 4 \\
\hline

B7 & Chatbot intelligent &
Développement et intégration du chatbot RAG (FAISS +
\texttt{all-MiniLM-L6-v2} + Groq) dans la plateforme. & Haute & 4 \\
\hline

B8 & Tests et déploiement &
Validation fonctionnelle, sécurisation, optimisation et
déploiement final de la plateforme. & Haute & 5 \\
\hline

\captionsetup{justification=centering,margin=2cm}
\caption{Product Backlog avec affectation des sprints}
\label{tab:product_backlog}
\end{longtable}

\vspace{0.4cm}

\begin{longtable}{|m{3cm}|m{4cm}|m{7cm}|}
\hline
\textbf{Release} & \textbf{Module} & \textbf{Sprint} \\
\hline
\endhead

Release 1 & Gestion des projets &
Sprint 1 : Enregistrement, suivi et paramétrage des projets \\
\hline

Release 1 & Gestion des utilisateurs &
Sprint 2 : Gestion des accès et pipeline ETL \\
\hline

Release 2 & Business Intelligence &
Sprint 3 : Tableaux de bord et reporting analytique \\
\hline

Release 2 & Intelligence Artificielle &
Sprint 4 : Prédiction des risques et chatbot RAG \\
\hline

Release 3 & Finalisation &
Sprint 5 : Sécurité, tests, optimisation et déploiement \\
\hline

\captionsetup{justification=centering,margin=2cm}
\caption{Répartition des releases et des sprints}
\label{tab:release_sprints}
\end{longtable}

\vspace{0.4cm}

\begin{longtable}{|m{1.8cm}|m{4.2cm}|m{8cm}|}
\hline
\textbf{Sprint} & \textbf{Période} & \textbf{Tâches principales} \\
\hline
\endhead

Sprint 1 & 05/01/2026 -- 31/01/2026 &
Analyse des besoins, modélisation initiale, développement du
module de gestion des projets PGD/PPP et du paramétrage
dynamique. \\
\hline

Sprint 2 & 01/02/2026 -- 28/02/2026 &
Gestion des utilisateurs, rôles et permissions ; développement
et test du pipeline ETL Python/Pandas. \\
\hline

Sprint 3 & 01/03/2026 -- 31/03/2026 &
Conception du schéma dimensionnel, développement des tableaux
de bord Superset et intégration Chart.js. \\
\hline

Sprint 4 & 01/04/2026 -- 30/04/2026 &
Entraînement du modèle XGBClassifier, développement du chatbot
RAG et intégration des deux modules dans FastAPI. \\
\hline

Sprint 5 & 01/05/2026 -- 01/06/2026 &
Tests fonctionnels et de performance, sécurisation des accès,
optimisation des requêtes et déploiement final. \\
\hline

\captionsetup{justification=centering,margin=2cm}
\caption{Calendrier de planification des sprints}
\label{tab:planning_sprints}
\end{longtable}

Cette organisation progressive garantit une livraison maîtrisée
des fonctionnalités et ménage la réactivité nécessaire pour
intégrer les retours des parties prenantes à l'issue de chaque
sprint.

% ============================================================
\section{Architecture générale}
% ============================================================

La figure~\ref{fig:architecture_generale} présente l'architecture
technique de la solution. Elle repose sur une organisation en cinq
couches fonctionnelles distinctes, qui favorisent la clarté des
responsabilités, la réutilisabilité des composants et la
maintenabilité de l'ensemble.

\begin{itemize}

    \item \textbf{Couche de présentation :} assurée par Joget,
    elle constitue le point d'entrée unique des utilisateurs.
    Elle prend en charge la saisie des données, le paramétrage
    dynamique des formulaires et le déclenchement des requêtes vers
    le backend via des appels API REST.

    \item \textbf{Couche de traitement :} orchestrée par FastAPI,
    elle centralise la réception des requêtes, l'appel au pipeline
    de traitement et l'invocation des modules d'IA. Les
    transformations sont réalisées en Python avec Pandas, qui
    applique les règles de nettoyage et de normalisation avant
    chargement dans l'entrepôt.

    \item \textbf{Couche de stockage :} les données se répartissent
    dans deux bases MySQL complémentaires :
    \begin{itemize}
        \item \texttt{jwdb} : base applicative générée
        automatiquement par Joget, hébergeant les données brutes
        issues des formulaires ainsi que l'entrepôt transformé,
        organisé selon le modèle dimensionnel de Kimball.
    \end{itemize}

    \item \textbf{Couche de visualisation :} l'entrepôt
    \texttt{cleaned\_dw} alimente Apache Superset pour les tableaux
    de bord décisionnels, complété par Chart.js pour les
    visualisations interactives intégrées dans Joget.

    \item \textbf{Modules d'intelligence artificielle :} deux
    modules enrichissent la chaîne décisionnelle :
    \begin{itemize}
        \item \textbf{Prédiction des risques :} un classificateur
        XGBClassifier évalue en temps réel le niveau de risque
        (faible, moyen, élevé, critique) et génère une explication
        contextuelle via \texttt{llama-3.3-70b-versatile} ;
        \item \textbf{Chatbot RAG :} combinant un index FAISS,
        le modèle \texttt{all-MiniLM-L6-v2} et l'API Groq, il
        fournit des réponses contextualisées aux questions des
        utilisateurs sur les projets environnementaux.
    \end{itemize}

\end{itemize}

\begin{figure}[H]
    \centering
    \includegraphics[width=0.95\textwidth]{img/archi.png}
    \caption{Architecture générale de la plateforme}
    \label{fig:architecture_generale}
\end{figure}

% ============================================================
\section{Environnement technique}
% ============================================================

Le tableau~\ref{tab:env_technique} recense les outils et technologies
mobilisés à chaque étape du projet, depuis la conception jusqu'au
déploiement. Trois critères ont guidé ces choix : la cohérence avec
l'architecture retenue, l'adéquation aux besoins fonctionnels et
la disponibilité en open source.

\begin{longtable}{|m{4cm}|m{4.5cm}|m{6cm}|}
\hline
\textbf{Phase} & \textbf{Outil} & \textbf{Rôle dans le projet} \\
\hline
\endhead

Développement applicatif (Low-Code)
& Joget DX
& Développement de l'interface utilisateur, des formulaires de saisie,
  du paramétrage dynamique et de la gestion des accès. \\
\hline

Développement backend
& FastAPI (Python)
& Exposition des endpoints REST, orchestration du pipeline de
  traitement, routage vers les modules IA et gestion des appels
  inter-services. \\
\hline

Traitement et pipeline ETL
& Python · Pandas
& Extraction des données brutes depuis \texttt{jwdb}, nettoyage,
  normalisation, enrichissement et chargement dans
  \texttt{cleaned\_dw}. \\
\hline

Stockage des données
& MySQL
& Hébergement de la base applicative Joget (\texttt{jwdb}) et de
  l'entrepôt décisionnel (\texttt{cleaned\_dw}). \\
\hline

Business Intelligence
& Apache Superset
& Création et publication des tableaux de bord interactifs connectés
  directement à \texttt{cleaned\_dw}. \\
\hline

Visualisation web
& Chart.js
& Intégration de graphiques dynamiques dans les interfaces Joget
  pour le suivi en temps réel des indicateurs projet. \\
\hline

Prédiction des risques (IA)
& XGBClassifier · scikit-learn · joblib
& Entraînement, évaluation et persistance du modèle de classification
  du risque projet sur 34 indicateurs. \\
\hline

Génération de langage naturel
& llama-3.3-70b-versatile (via API Groq)
& Production des explications contextuelles des prédictions de risque
  et génération des réponses du chatbot. \\
\hline

Chatbot RAG
& FAISS · SentenceTransformers (\texttt{all-MiniLM-L6-v2})
& Construction de l'index vectoriel et recherche sémantique dans la
  base documentaire pour alimenter le chatbot en contexte pertinent. \\
\hline

Gestion des dépendances
& pip · requirements.txt
& Gestion et reproductibilité de l'environnement Python du backend. \\
\hline

Versionnement du code
& Git · GitHub
& Suivi des versions du code source, collaboration et sauvegarde du
  dépôt du projet. \\
\hline

\captionsetup{justification=centering,margin=2cm}
\caption{Environnement technique du projet}
\label{tab:env_technique}
\end{longtable}

L'ensemble forme un écosystème cohérent et majoritairement open
source, qui évite toute dépendance à des licences propriétaires
coûteuses tout en offrant les performances attendues d'un système
décisionnel orienté données et intelligence artificielle.

% ============================================================
\section{Inspection des données source}
% ============================================================

Les données exploitées dans ce projet proviennent de deux niveaux
distincts de la chaîne de traitement.

Le premier niveau regroupe les \textbf{tables brutes} générées
automatiquement par Joget lors de l'import historique des données
depuis les fichiers Excel fournis par Wevioo :
\texttt{app\_fd\_base\_pgd}, \texttt{app\_fd\_base\_pgd\_etude}
et \texttt{app\_fd\_base\_ppp}. Elles constituent la source
primaire du pipeline ETL.

Le second niveau correspond aux \textbf{tables nettoyées} —
\texttt{app\_fd\_pgd\_clean}, \texttt{app\_fd\_pgd\_etude\_clean}
et \texttt{app\_fd\_ppp\_clean} — produites par le pipeline ETL
Python développé dans le cadre de ce projet. C'est sur elles que
s'appuient l'entrepôt de données et les modules d'IA.

L'inspection porte sur les tables brutes sources. Pour chaque
colonne, deux dimensions sont évaluées : la \textbf{qualité}
(cohérence, complétude) et l'\textbf{utilité analytique}
(pertinence pour l'entrepôt et les tableaux de bord).
Le symbole \checkmark\ indique une valeur positive, \texttimes\
une anomalie ou une exclusion justifiée.

% ============================================================
\subsection{Inspection de la table \texttt{app\_fd\_base\_pgd}}
% ============================================================

Cette table regroupe les projets opérationnels de gestion des
déchets importés depuis les fichiers historiques Excel. Elle
compte \textbf{16 colonnes}.

\begin{longtable}{
  | >{\RaggedRight}p{3cm}
  | >{\RaggedRight}p{3.2cm}
  | >{\centering\arraybackslash}p{1.2cm}
  | >{\centering\arraybackslash}p{0.9cm}
  | >{\centering\arraybackslash}p{0.8cm}
  | >{\RaggedRight}p{3.8cm} |
}
\hline
\textbf{Nom} & \textbf{Description} & \textbf{Type} &
\textbf{Qual.} & \textbf{Util.} & \textbf{Commentaire} \\
\hline
\endhead

id & Identifiant unique (UUID Joget) & String &
\checkmark & \checkmark & Clé primaire \\
\hline
dateCreated & Date de création & String &
\checkmark & \checkmark & Conversion datetime requise \\
\hline
dateModified & Date de modification & String &
\checkmark & \checkmark & Conversion datetime requise \\
\hline
createdBy & Login créateur & String &
\checkmark & \texttimes & Non analytique — exclu du DW \\
\hline
createdByName & Nom créateur & String &
\checkmark & \texttimes & Non analytique — exclu du DW \\
\hline
modifiedBy & Login modificateur & String &
\checkmark & \texttimes & Non analytique — exclu du DW \\
\hline
modifiedByName & Nom modificateur & String &
\checkmark & \texttimes & Non analytique — exclu du DW \\
\hline
c\_projet & Intitulé du projet & String &
\checkmark & \checkmark & \\
\hline
c\_promoteur & Organisme promoteur & String &
\checkmark & \checkmark & \\
\hline
c\_activite & Activités réalisées & String &
\checkmark & \checkmark & Texte long — usage descriptif \\
\hline
c\_objectifs & Objectifs du projet & String &
\checkmark & \checkmark & Texte long — usage descriptif \\
\hline
c\_context & Contexte du projet & String &
\checkmark & \checkmark & Texte long — usage descriptif \\
\hline
c\_bailleur & Bailleur de fonds & String &
\checkmark & \checkmark & \\
\hline
c\_annee & Année du projet & Float64 &
\checkmark & \checkmark & Convertir en Int64 dans le DW \\
\hline
c\_region & Région géographique & String &
\checkmark & \checkmark & \\
\hline
c\_budget & Budget alloué & String &
\texttimes & \checkmark & Format texte libre (ex. : ``9 millions de DT'') — extraction numérique et conversion TND effectuées lors du nettoyage \\
\hline

\caption{Inspection de la table brute \texttt{app\_fd\_base\_pgd}}
\label{tab:inspection_pgd}
\end{longtable}

% ============================================================
\subsection{Inspection de la table \texttt{app\_fd\_base\_pgd\_etude}}
% ============================================================

Cette table regroupe les études et rapports associés aux projets
PGD. Elle contient \textbf{20 colonnes} et inclut notamment
des liens PDF et des références vers les tables de paramétrage.

\begin{longtable}{
  | >{\RaggedRight}p{3cm}
  | >{\RaggedRight}p{3.2cm}
  | >{\centering\arraybackslash}p{1.2cm}
  | >{\centering\arraybackslash}p{0.9cm}
  | >{\centering\arraybackslash}p{0.8cm}
  | >{\RaggedRight}p{3.8cm} |
}
\hline
\textbf{Nom} & \textbf{Description} & \textbf{Type} &
\textbf{Qual.} & \textbf{Util.} & \textbf{Commentaire} \\
\hline
\endhead

id & Identifiant unique (UUID Joget) & String &
\checkmark & \checkmark & Clé primaire \\
\hline
dateCreated & Date de création & String &
\checkmark & \checkmark & Conversion datetime requise \\
\hline
dateModified & Date de modification & String &
\checkmark & \checkmark & Conversion datetime requise \\
\hline
createdBy & Login créateur & String &
\checkmark & \texttimes & Non analytique — exclu du DW \\
\hline
createdByName & Nom créateur & String &
\checkmark & \texttimes & Non analytique — exclu du DW \\
\hline
modifiedBy & Login modificateur & String &
\checkmark & \texttimes & Non analytique — exclu du DW \\
\hline
modifiedByName & Nom modificateur & String &
\checkmark & \texttimes & Non analytique — exclu du DW \\
\hline
c\_intitule & Intitulé de l'étude & String &
\checkmark & \checkmark & \\
\hline
c\_expert & Expert responsable & String &
\checkmark & \checkmark & \\
\hline
c\_annee & Année de l'étude & Int64 &
\checkmark & \checkmark & \\
\hline
c\_cout & Coût de l'étude (DT) & Int64 &
\checkmark & \checkmark & Indicateur financier clé \\
\hline
c\_domaine & Domaine de l'étude & String &
\checkmark & \checkmark & FK → param\_domaine\_pgd \\
\hline
c\_theme & Thème de l'étude & String &
\checkmark & \checkmark & FK → param\_theme\_pgd \\
\hline
c\_categorie & Catégorie de l'étude & String &
\checkmark & \checkmark & FK → param\_categorie\_pgd \\
\hline
c\_bailleur & Bailleur de fonds & String &
\checkmark & \checkmark & FK → param\_bailleur\_pgd \\
\hline
c\_beneficiaires & Bénéficiaires & String &
\checkmark & \checkmark & Texte libre \\
\hline
c\_acteur\_institutionnel & Acteur institutionnel & String &
\checkmark & \checkmark & \\
\hline
c\_partenaire\_implementation & Partenaire d'implémentation & String &
\checkmark & \checkmark & \\
\hline
c\_cadre & Cadre programmatique & String &
\checkmark & \checkmark & \\
\hline
c\_lien\_pdf & Lien URL vers le PDF & String &
\checkmark & \checkmark & Utile pour la bibliothèque \\
\hline

\caption{Inspection de la table brute \texttt{app\_fd\_base\_pgd\_etude}}
\label{tab:inspection_pgd_etude}
\end{longtable}

% ============================================================
\subsection{Inspection de la table \texttt{app\_fd\_base\_ppp}}
% ============================================================

Avec \textbf{73 colonnes}, c'est la table la plus volumineuse du
corpus. Deux particularités structurelles ressortent de son
inspection : les dates sont encodées au format sériel Excel
(entier = nombre de jours depuis le 01/01/1900), ce qui impose
une conversion systématique dans le pipeline ETL ; plusieurs
colonnes contiennent des UUID pointant vers les tables de
paramétrage Joget (clés étrangères). Par souci de lisibilité,
les colonnes FK n'indiquent que le suffixe de la table cible.

\begin{longtable}{
  | >{\RaggedRight}p{3cm}
  | >{\RaggedRight}p{3.2cm}
  | >{\centering\arraybackslash}p{1.2cm}
  | >{\centering\arraybackslash}p{0.9cm}
  | >{\centering\arraybackslash}p{0.8cm}
  | >{\RaggedRight}p{3.8cm} |
}
\hline
\textbf{Nom} & \textbf{Description} & \textbf{Type} &
\textbf{Qual.} & \textbf{Util.} & \textbf{Commentaire} \\
\hline
\endhead

id & Identifiant unique & String &
\checkmark & \checkmark & Clé primaire \\
\hline
dateCreated & Date de création & String &
\checkmark & \checkmark & Conversion datetime requise \\
\hline
dateModified & Date de modification & String &
\checkmark & \checkmark & Conversion datetime requise \\
\hline
createdBy & Login créateur & String &
\checkmark & \texttimes & Non analytique — exclu du DW \\
\hline
createdByName & Nom créateur & String &
\checkmark & \texttimes & Non analytique — exclu du DW \\
\hline
modifiedBy & Login modificateur & String &
\checkmark & \texttimes & Non analytique — exclu du DW \\
\hline
modifiedByName & Nom modificateur & String &
\checkmark & \texttimes & Non analytique — exclu du DW \\
\hline
c\_nom\_projet & Intitulé du projet PPP & String &
\checkmark & \checkmark & \\
\hline
c\_numero & Numéro du projet & Int64 &
\checkmark & \checkmark & \\
\hline
c\_cadre & Cadre budgétaire (ex. LF2020) & String &
\checkmark & \checkmark & \\
\hline
c\_odd & ODD associé au projet & String &
\checkmark & \checkmark & Texte libre \\
\hline
c\_gouvernorat & Gouvernorat & Int64 &
\checkmark & \checkmark & FK → app\_fd\_gouvernorat \\
\hline
c\_delegation & Délégation & Int64 &
\checkmark & \checkmark & FK → app\_fd\_delegation \\
\hline
c\_commune & Commune & Int64 &
\checkmark & \checkmark & Clé numérique \\
\hline
c\_localite & Localité précise & String &
\texttimes & \checkmark & 50\,\% de valeurs manquantes \\
\hline
c\_axe & Axe stratégique & UUID &
\checkmark & \checkmark & FK → param\_axe\_ppp \\
\hline
c\_sous\_axe & Sous-axe stratégique & UUID &
\checkmark & \checkmark & FK → param\_sous\_axe\_ppp \\
\hline
c\_classification & Classification du projet & UUID &
\checkmark & \checkmark & FK → param\_classification \\
\hline
c\_sous\_classification & Sous-classification & UUID &
\checkmark & \checkmark & FK → param\_sous\_class\_ppp \\
\hline
c\_initiateur & Initiateur du projet & UUID &
\checkmark & \checkmark & FK → param\_initiateur\_ppp \\
\hline
c\_owner & Structure porteuse & String &
\checkmark & \checkmark & Valeur textuelle directe \\
\hline
c\_acteur\_dimplementation & Acteur d'implémentation & UUID &
\checkmark & \checkmark & FK → param\_act\_imp\_ppp \\
\hline
c\_delimination & Délimitation géographique & UUID &
\checkmark & \checkmark & FK → param\_delim\_geo\_ppp \\
\hline
c\_mode\_financement & Mode de financement & UUID &
\checkmark & \checkmark & FK → param\_mode\_finance \\
\hline
c\_mode\_passation\_marche & Mode de passation marché & UUID &
\checkmark & \checkmark & FK → param\_pass\_march \\
\hline
c\_tutelle\_sectorielle & Tutelle sectorielle & UUID &
\checkmark & \checkmark & FK → param\_tut\_sec\_ppp \\
\hline
c\_zonage & Zonage du projet & UUID &
\checkmark & \checkmark & FK → param\_zonage\_ppp \\
\hline
c\_donateur & Donateur & UUID &
\texttimes & \checkmark & 80\,\% de valeurs manquantes ; FK → param\_donateur \\
\hline
c\_bailleur\_fonds & Bailleur de fonds & String &
\texttimes & \checkmark & 70\,\% de valeurs manquantes \\
\hline
c\_situation\_actuelle\_projet & Situation actuelle & String &
\checkmark & \checkmark & \\
\hline
c\_phase\_en\_cours & Phase opérationnelle & String &
\checkmark & \checkmark & \\
\hline
c\_exercice\_budgetaire & Exercice budgétaire & Int64 &
\checkmark & \checkmark & \\
\hline
c\_budget\_global\_planifie & Budget planifié (DT) & Int64 &
\checkmark & \checkmark & Indicateur financier clé \\
\hline
c\_budget\_global\_consomme & Budget consommé (DT) & Float64 &
\checkmark & \checkmark & Indicateur financier clé \\
\hline
c\_solde\_budget\_restant & Solde restant (DT) & Float64 &
\checkmark & \checkmark & Planifié − Consommé \\
\hline
c\_budget\_realisation\_planifie & Budget réalisation planifié & Int64 &
\checkmark & \checkmark & \\
\hline
c\_budget\_realisation\_consomme & Budget réalisation consommé & Float64 &
\checkmark & \checkmark & \\
\hline
c\_budget\_etude\_planifie & Budget études planifié & Int64 &
\checkmark & \checkmark & \\
\hline
c\_budget\_etudes\_consomme & Budget études consommé & Float64 &
\checkmark & \checkmark & \\
\hline
c\_budget\_controle\_planifie & Budget contrôle planifié & Int64 &
\checkmark & \checkmark & \\
\hline
c\_budget\_controle\_consomme & Budget contrôle consommé & Float64 &
\checkmark & \checkmark & \\
\hline
c\_cout\_foncier & Coût foncier (DT) & Int64 &
\checkmark & \checkmark & \\
\hline
c\_part\_budget\_etat & Part budget État (DT) & Int64 &
\checkmark & \checkmark & \\
\hline
c\_part\_credits & Part crédits (DT) & Int64 &
\checkmark & \checkmark & \\
\hline
c\_part\_dons & Part dons (DT) & Int64 &
\checkmark & \checkmark & \\
\hline
c\_taux\_avancement\_institutionnel & Avancement institutionnel & Float64 &
\checkmark & \checkmark & Valeur décimale 0–1 \\
\hline
c\_taux\_avancement\_passation & Avancement passation marchés & Float64 &
\checkmark & \checkmark & Valeur décimale 0–1 \\
\hline
c\_nombre\_marches & Nombre de marchés & Int64 &
\checkmark & \checkmark & \\
\hline
c\_retard\_avance\_en\_jours & Retard / avance (jours) & Int64 &
\checkmark & \checkmark & Négatif = avance, positif = retard \\
\hline
c\_date\_debut\_prevue & Date début prévue & Int64 &
\texttimes & \checkmark & Format sériel Excel — conversion obligatoire \\
\hline
c\_date\_fin\_prevue & Date fin prévue & Int64 &
\texttimes & \checkmark & Format sériel Excel \\
\hline
c\_date\_debut\_reelle & Date début réelle & Int64 &
\texttimes & \checkmark & Format sériel Excel \\
\hline
c\_date\_fin\_reelle & Date fin réelle & Int64 &
\texttimes & \checkmark & Format sériel Excel \\
\hline
c\_date\_approbation\_cnapp & Date approbation CNAPP & Int64 &
\texttimes & \checkmark & Format sériel Excel \\
\hline
c\_date\_etude\_faisabilite & Date étude faisabilité & Int64 &
\texttimes & \checkmark & Format sériel Excel \\
\hline
c\_nombre\_beneficiaires & Nombre de bénéficiaires & Int64 &
\checkmark & \checkmark & \\
\hline
c\_emplois\_directs\_crees & Emplois directs créés & Int64 &
\checkmark & \checkmark & \\
\hline
c\_part\_femmes\_execution & Part femmes — exécution & Float64 &
\checkmark & \checkmark & Valeur 0–1 \\
\hline
c\_part\_femmes\_emplois & Part femmes — emplois créés & Float64 &
\checkmark & \checkmark & Valeur 0–1 \\
\hline
c\_part\_femmes\_beneficiaires & Part femmes — bénéficiaires & Float64 &
\checkmark & \checkmark & Valeur 0–1 \\
\hline
c\_incidents\_corporels & Incidents corporels & Int64 &
\checkmark & \checkmark & Indicateur sécurité \\
\hline
c\_penalites\_amendes & Pénalités et amendes (DT) & Int64 &
\checkmark & \checkmark & \\
\hline
c\_etude\_impact\_environnement & Étude impact environnemental & String &
\checkmark & \checkmark & Réalisée / Non réalisée \\
\hline
c\_normes\_techniques & Normes techniques applicables & String &
\checkmark & \checkmark & Texte libre \\
\hline
c\_contractant\_prive & Nom du contractant privé & String &
\checkmark & \checkmark & \\
\hline
c\_rne\_contractant\_prive & RNE contractant privé & String &
\checkmark & \texttimes & Administratif — non analytique \\
\hline
c\_taille\_contractant\_prive & Taille contractant (PME/Grande) & String &
\checkmark & \checkmark & \\
\hline
c\_responsable\_institutionnel & Responsable institutionnel & String &
\checkmark & \checkmark & \\
\hline
c\_responsable\_suivi & Responsable de suivi & String &
\checkmark & \checkmark & \\
\hline
c\_projet\_parent & Projet parent (si existant) & String &
\checkmark & \checkmark & 70\,\% de valeurs manquantes \\
\hline
c\_projet\_child & Projet enfant & Float64 &
\texttimes & \texttimes & 100\,\% de valeurs manquantes — colonne supprimée \\
\hline
c\_commune\_s8 & Commune section 8 & Float64 &
\texttimes & \texttimes & 100\,\% de valeurs manquantes — colonne supprimée \\
\hline

\captionsetup{justification=centering,margin=2cm}
\caption{Inspection de la table brute \texttt{app\_fd\_base\_ppp}}
\label{tab:inspection_ppp}
\end{longtable}

% ============================================================
\subsection{Identification des tables de faits et de dimensions}
% ============================================================

L'entrepôt \texttt{cleaned\_dw} s'organise autour d'un schéma
dimensionnel composé de trois tables de faits et de vingt et une
tables de dimensions.

\subsubsection{Tables de faits}

Les trois tables de faits sont le produit direct du pipeline ETL.

\begin{longtable}{|p{4.5cm}|p{4cm}|p{6cm}|}
\hline
\textbf{Table brute (source)} & \textbf{Table nettoyée (DW)} &
\textbf{Description} \\
\hline
\endhead

\texttt{app\_fd\_base\_pgd} &
\texttt{app\_fd\_pgd\_clean} &
Projets opérationnels de gestion des déchets — budget normalisé en
TND, année convertie en entier. \\
\hline

\texttt{app\_fd\_base\_pgd\_etude} &
\texttt{app\_fd\_pgd\_etude\_clean} &
Études et rapports PGD — coût normalisé, colonnes de traçabilité
exclues. \\
\hline

\texttt{app\_fd\_base\_ppp} &
\texttt{app\_fd\_ppp\_clean} &
Projets PPP — dates converties, indicateurs de risque calculés
(\texttt{risque\_retard}, \texttt{risque\_depassement}),
colonnes vides supprimées. \\
\hline

\captionsetup{justification=centering,margin=2cm}
\caption{Correspondance tables brutes / tables nettoyées}
\label{tab:tables_faits}
\end{longtable}

\subsubsection{Tables de dimensions}

Les tables de dimensions proviennent des modules de paramétrage
Joget. Elles portent un identifiant unique — stocké comme clé
étrangère dans les tables de faits — et un libellé décrivant
l'entité référencée.

\begin{longtable}{|p{5.5cm}|p{9cm}|}
\hline
\textbf{Table de dimension} & \textbf{Description} \\
\hline
\endhead

\texttt{app\_fd\_param\_domaine\_pgd} &
Domaines des études PGD (ex. : Environnement, Économie circulaire).
Référencé par \texttt{c\_domaine}. \\
\hline
\texttt{app\_fd\_param\_theme\_pgd} &
Thèmes des études PGD. Référencé par \texttt{c\_theme}. \\
\hline
\texttt{app\_fd\_param\_categorie\_pgd} &
Catégories des études PGD. Référencé par \texttt{c\_categorie}. \\
\hline
\texttt{app\_fd\_param\_bailleur\_pgd} &
Bailleurs de fonds PGD. Référencé par \texttt{c\_bailleur}. \\
\hline
\texttt{app\_fd\_param\_axe\_ppp} &
Axes stratégiques PPP. Référencé par \texttt{c\_axe}. \\
\hline
\texttt{app\_fd\_param\_sous\_axe\_ppp} &
Sous-axes stratégiques. Référencé par \texttt{c\_sous\_axe}. \\
\hline
\texttt{app\_fd\_param\_classification} &
Classifications des projets PPP. Référencé par \texttt{c\_classification}. \\
\hline
\texttt{app\_fd\_param\_sous\_class\_ppp} &
Sous-classifications PPP. Référencé par \texttt{c\_sous\_classification}. \\
\hline
\texttt{app\_fd\_param\_act\_imp\_ppp} &
Acteurs d'implémentation. Référencé par \texttt{c\_acteur\_dimplementation}. \\
\hline
\texttt{app\_fd\_param\_delim\_geo\_ppp} &
Délimitations géographiques. Référencé par \texttt{c\_delimination}. \\
\hline
\texttt{app\_fd\_param\_initiateur\_ppp} &
Initiateurs des projets. Référencé par \texttt{c\_initiateur}. \\
\hline
\texttt{app\_fd\_param\_owner\_ppp} &
Structures porteuses. Référencé par \texttt{c\_owner}. \\
\hline
\texttt{app\_fd\_param\_mode\_finance} &
Modes de financement. Référencé par \texttt{c\_mode\_financement}. \\
\hline
\texttt{app\_fd\_param\_pass\_march} &
Modes de passation de marché. Référencé par \texttt{c\_mode\_passation\_marche}. \\
\hline
\texttt{app\_fd\_param\_tut\_sec\_ppp} &
Tutelles sectorielles. Référencé par \texttt{c\_tutelle\_sectorielle}. \\
\hline
\texttt{app\_fd\_param\_zonage\_ppp} &
Zones géographiques. Référencé par \texttt{c\_zonage}. \\
\hline
\texttt{app\_fd\_param\_situation\_pj} &
Situations des projets (En cours, Bloqué, Clôturé). \\
\hline
\texttt{app\_fd\_param\_donateur} &
Donateurs associés aux projets PPP (80\,\% de valeurs nulles). \\
\hline
\texttt{app\_fd\_gouvernorat} &
Référentiel des gouvernorats tunisiens. Clé entière. \\
\hline
\texttt{app\_fd\_delegation} &
Référentiel des délégations. Clé entière. \\
\hline
\texttt{app\_fd\_secteur} &
Référentiel des secteurs pour la granularité infranationale. \\
\hline

\captionsetup{justification=centering,margin=2cm}
\caption{Tables de dimensions de l'entrepôt de données}
\label{tab:tables_dimensions}
\end{longtable}

% ============================================================
\section{Conception du modèle physique de données}
% ============================================================

L'inspection des tables sources et l'identification des entités
dimensionnelles ont conduit à la conception du modèle physique de
l'entrepôt \texttt{cleaned\_dw}. Ce modèle est le socle sur lequel
reposent les traitements analytiques, les tableaux de bord et les
modules d'IA.

\subsection{Type de schéma retenu}

L'entrepôt ne relève pas d'un simple schéma en étoile. La
coexistence de trois tables de faits distinctes partageant des
dimensions communes — notamment \texttt{app\_fd\_param\_bailleur\_pgd},
\texttt{app\_fd\_gouvernorat} et \texttt{app\_fd\_delegation} —
en fait un \textbf{schéma en constellation} (\textit{galaxy schema}).
Ce choix évite la redondance des référentiels et garantit la
cohérence des données entre les trois domaines analytiques couverts.

\subsection{Organisation des tables}

Le modèle physique s'articule autour de trois pôles fonctionnels.

\begin{itemize}
    \item \textbf{Pôle PGD opérationnel :} la table
    \texttt{app\_fd\_pgd\_clean} centralise les projets de gestion
    des déchets avec budget normalisé en TND et année de réalisation.
    Elle partage la dimension bailleur avec le pôle suivant.

    \item \textbf{Pôle PGD études :} \texttt{app\_fd\_pgd\_etude\_clean}
    héberge les études et rapports PGD. Elle partage la dimension
    bailleur et dispose de trois dimensions propres : domaine, thème
    et catégorie.

    \item \textbf{Pôle PPP :} \texttt{app\_fd\_ppp\_clean} est la
    table la plus dense — dix-neuf clés étrangères couvrant axes,
    classifications, financements, acteurs, indicateurs de performance
    et localisation. Elle intègre deux colonnes calculées par le
    pipeline ETL — \texttt{risque\_retard} et
    \texttt{risque\_depassement} — directement utilisées comme
    features d'entrée du modèle XGBClassifier.
\end{itemize}

\subsection{Hiérarchie géographique}

La modélisation géographique mérite une attention particulière.
\texttt{app\_fd\_delegation} référence \texttt{app\_fd\_gouvernorat}
par clé étrangère entière, créant une hiérarchie administrative en
trois niveaux : gouvernorat → délégation → secteur. Cette structure
permet à Superset d'agréger les données à différentes granularités
territoriales.

\subsection{Représentation du modèle physique}

\begin{figure}[H]
    \centering
    \includegraphics[width=1.1\textwidth]{img/constellation_schema_dw.png}
    \caption{Modèle physique de l'entrepôt de données
             — schéma en constellation}
    \label{fig:modele_physique}
\end{figure}

Conçu selon les principes de la modélisation dimensionnelle de
Kimball, ce schéma offre des performances de requêtage optimales
pour Superset tout en restant suffisamment lisible pour les
développeurs chargés de la maintenance du pipeline.

\newpage

% ============================================================
\section*{Conclusion}
% ============================================================

Ce chapitre a conduit une démarche d'analyse structurée, des besoins
métiers jusqu'au modèle physique de données.

Les sept axes de besoins d'affaires ont fourni le cadre de référence
pour tous les arbitrages de conception. Deux profils d'utilisateurs
et huit domaines fonctionnels ont été identifiés et formalisés. Les
exigences non fonctionnelles — fiabilité, sécurité, performance,
scalabilité — complètent ce tableau en définissant les contraintes
de qualité que la plateforme doit respecter en toutes circonstances.

La planification en cinq sprints Agile garantit une progression
maîtrisée et réversible des développements. L'architecture en cinq
couches et l'environnement technique majoritairement open source
traduisent concrètement ces besoins en choix de conception.
Enfin, l'inspection des données sources a mis en évidence les
principales anomalies à corriger et a conduit à la définition du
schéma en constellation de l'entrepôt, composé de trois tables de
faits et de vingt et une dimensions.

Ces fondations posées, le chapitre suivant entre dans la réalisation
technique : pipeline ETL, développement des interfaces et intégration
des modules d'intelligence artificielle.